\section{Proposed Empirical Design}
\label{sec:design}

This section presents the specifications, defines the treatment, describes the data, and pairs each identifying assumption with its threat and the specification that would probe it.

\subsection{Specification}

The baseline continuous-treatment difference-in-differences, estimated on the North African sample alone, is:
\begin{equation}
\label{eq:baseline}
\ln M_{pct} = \beta \, (\tau_p \times C_{st}) + \phi_{pc} + \psi_{ct} + u_{pct},
\end{equation}
where $M_{pct}$ is the import value (EUR) of product $p$ from partner $c$ in year $t$; $\tau_p$ is product-level carbon-cost exposure; $C_{st}$ is the effective carbon cost of the supplying sector $s = s(p)$; $\phi_{pc}$ are product-by-partner fixed effects; and $\psi_{ct}$ are partner-by-year fixed effects. The preferred triple-difference specification pools the corridor with a comparison partner group:
\begin{equation}
\label{eq:triple}
\ln M_{pct} = \beta \, (\tau_p \times C_{st} \times L_c) + \phi_{pc} + \psi_{ct} + \eta_{pt} + u_{pct},
\end{equation}
where $L_c$ indicates the lower-regulation corridor (Morocco, Algeria, Egypt, Tunisia) and $\eta_{pt}$ are product-by-year fixed effects. A parallel specification would replace $\ln M_{pct}$ with quantity to separate volume from unit-price movements.

\subsection{Variable Definitions}

\paragraph{Carbon-cost exposure ($\tau_p$).}
The CO$_2$ emission intensity of the sector supplying product $p$ under the European industrial activity classification NACE Rev.~2, mapped to products through the RAMON correspondence between the Combined Nomenclature (CN) and NACE \citep{ramon}, fixed at its 2005--2007 average to prevent endogenous intensity responses. The transparent alternative is a binary CBAM Annex~I indicator \citep{cbam2023}.

\paragraph{Effective carbon cost ($C_{st}$).}
\begin{equation}
\label{eq:cst}
C_{st} = P_t \times (1 - a_{st}),
\end{equation}
where $P_t$ is the annual average EUA price \citep{ember} and $a_{st}$ is the sector-year free-allocation share from the EU Transaction Log \citep{eutl}. Free-allocation attenuation is part of the treatment definition: a sector with 100 per cent free allocation faces zero effective cost regardless of the EUA price. An objection is that tradable allowances carry the full permit price as marginal opportunity cost, so lump-sum allocation would leave the output margin unaffected. EU allocation is not lump-sum in this strict sense: activity-level thresholds, closure provisions, and partial-cessation rules tie allowances to continued production, which gives free allocation genuine marginal bite \citep{ellerman2016}. A second concern is mechanical endogeneity: contemporaneous emissions enter the denominator of the allocation share, so demand shocks that cut emissions would raise $a_{st}$ by construction. The baseline therefore computes $a_{st}$ relative to fixed 2005--2007 verified emissions. In Phases~I and~II, $a_{st}$ can exceed one; the baseline caps $a_{st}$ at one. The contemporaneous share and the uncapped version are retained as robustness checks.

\paragraph{Comparison group ($L_c = 0$).}
The choice of comparison group determines what $\beta$ identifies. Three candidates merit consideration. Turkey participates in an EU customs union that eliminates industrial tariffs, so carbon-cost effects are difficult to separate from pre-existing preferential-access effects; Turkey has also begun developing its own emissions trading system, which introduces partial regulatory convergence towards the end of the sample period. The Western Balkans (Albania, Bosnia and Herzegovina, Kosovo, Montenegro, North Macedonia, Serbia) share geographic proximity, operate under Stabilisation and Association Agreements with comparable tariff schedules, and lacked carbon-pricing regimes throughout the sample period, a cleaner regulatory comparison; their limitation is smaller export volumes in emission-intensive goods, which introduces noise. An intra-EU sourcing benchmark would control for EU consumption shocks but faces EU~ETS regulation by construction, so it is not a lower-regulation comparator in the standard sense. The proposed design would report results with each group and discuss sensitivity to this choice.

\paragraph{Fixed-effect absorption.}
In Equation~\ref{eq:triple}, $\eta_{pt}$ absorbs $\tau_p \times C_{st}$, the product-level effect regardless of partner; $\psi_{ct}$ absorbs $C_{st} \times L_c$ and all partner-level macro and exchange-rate shocks; $\phi_{pc}$ absorbs $\tau_p \times L_c$. Only the triple interaction survives. Identification comes from the differential response of high-exposure products from the lower-regulation corridor, purged of global product-specific shocks and partner-specific conditions.

\subsection{Sample and Data}

The sample covers 2005 to 2021, spanning Phase~I through the first year of Phase~IV. Pre-2005 years (2000--2004) are reserved for pre-trend validation; the Directive dates from 2003, so 2000 to 2002 constitutes a clean pre-announcement period and 2003 to 2004 a potential anticipation window. The window ends in 2021 by design: the 2022 energy-price shock, driven by the Russian invasion of Ukraine, distorted both EUA prices and bilateral trade flows in ways unrelated to the leakage channel.

Bilateral trade data at the CN8 level would come from Eurostat Comext (dataset DS-045409; \citealp{comext}), which provides value and quantity by reporting Member State, partner country, product code, and year. Comext coverage of North African partners is comprehensive for goods trade. UN Comtrade data at the six-digit Harmonised System level (HS6) would be the fallback for missing partner-product cells \citep{comtrade}; the coarser classification reduces product-level variation but has wider country coverage. Sectoral CO$_2$ emission intensities would be computed from Eurostat air emissions accounts, defined as tonnes of CO$_2$ per unit of gross output in the corresponding NACE Rev.~2 sector, and mapped to trade products through the RAMON CN-to-NACE correspondence \citep{ramon}.

\subsection{Identifying Assumptions and Threats}

\begin{enumerate}

\item \emph{Parallel trends.} Absent the EU~ETS, differential trends in high-exposure imports from the corridor (relative to comparison partners) would match those for low-exposure products. An event-study specification interacting $\tau_p$ with year indicators (base 2004) would inspect 2000--2002 for pre-trends and 2003--2004 for anticipation.

\item \emph{No confounding product shocks.} Global shocks correlated with emission intensity do not differentially affect the corridor versus comparison partners. $\eta_{pt}$ absorbs all product-by-year variation; a remaining violation would require a shock simultaneously product-specific, intensity-correlated, and corridor-differential. In the baseline Equation~\ref{eq:baseline}, product-level commodity-price controls partially address this.

\item \emph{Carbon-price exogeneity.} The EUA price, set by EU-wide policy and macro factors, is plausibly exogenous to individual bilateral flows. $\psi_{ct}$ absorbs partner-year variation, so identification comes from within-partner-year, across-product variation.

\item \emph{No coincident trade-policy changes.} Euro-Mediterranean association agreement tariff phase-outs could confound the estimate if they coincided with EU~ETS phases and disproportionately affected emission-intensive goods. The agreements entered into force at different dates (Tunisia 1998, Morocco 2000, Egypt 2004, Algeria 2005), and industrial tariff phase-outs were largely completed before the post-2018 price acceleration. The triple-difference mitigates this insofar as comparison partners face similar preferential arrangements.

\item \emph{Over-allocation attenuation.} $a_{st}$ exceeding one in Phases~I and~II mechanically weakens the treatment. Capping $a_{st}$ at one sets $C_{st} = 0$ for over-allocated sector-years; the uncapped version is retained as a robustness check.

\item \emph{CBAM anticipation.} The July 2021 proposal could alter 2021 trade patterns, but anticipation of a border charge deters leakage and biases $\beta$ toward zero. The event-study specification would detect any anomalous 2021 coefficient.

\item \emph{Pandemic disruption.} The years 2020 and 2021 fall within the sample and coincide with severe trade and supply-chain disruption. $\psi_{ct}$ absorbs partner-specific pandemic severity and $\eta_{pt}$ absorbs product-specific demand collapses; a residual violation would require pandemic effects that were simultaneously corridor-specific and correlated with emission intensity. A robustness check would re-estimate the specification excluding 2020 and 2021.

\item \emph{Corridor-specific energy costs.} Algeria and Egypt produce natural gas and regulate domestic energy prices, so the corridor's cost advantage in gas-intensive products widens when European gas prices rise; European gas prices in turn co-move with the EUA price, most strongly after 2018. This confound is corridor-specific, intensity-correlated, and time-varying, so no fixed effect in Equation~\ref{eq:triple} absorbs it. The specification would add a control interacting $\tau_p \times L_c$ with the EU--partner gas-price differential; without this control, $\beta$ would bundle the carbon-cost asymmetry with the energy-cost asymmetry.

\end{enumerate}

\subsection{Inference and Implementation}

Standard errors would be clustered two-way by product and partner-year \citep{cameron2011}. The small number of partner clusters raises finite-cluster concerns; sector-level clustering and a wild cluster bootstrap would address this. The design would be implemented with ordinary least squares (OLS) and high-dimensional fixed effects (\texttt{fixest}; \citealp{berge2021}) and with Poisson pseudo-maximum-likelihood (PPML) estimation (\texttt{fepois}) to handle zero trade flows consistently \citep{santos2006}.

\subsection{Validation}

Two placebo exercises would assess specificity: estimation on non-tradable products for which the import channel is inoperative, and estimation on low-intensity products (below the 25th percentile of $\tau_p$). A heterogeneity analysis would estimate $\beta$ separately for the main CBAM product groups, whose trade exposure differs: cement's low value-to-weight ratio confines trade to short sea routes, which the Mediterranean corridor provides, whereas steel and fertilisers trade globally. Concentration of the effect in products for which corridor trade is feasible would support the leakage interpretation. A positive and significant $\beta$ would be consistent with corridor-specific leakage; a precisely estimated zero would suggest free allocation has been sufficient; an imprecise zero would leave the question unresolved and motivate more disaggregated data.
